% =============================================================================
% 08_appendix.tex — lane WC
% Sources: .omo/evidence/task-F3-tensixflow.log (F3 final verification wave),
%          docs/hardware_path.md, docs/performance_model.md,
%          kernels/handwritten_baselines/*, .omo/plans glossary terms
% =============================================================================
\section{Appendix}\label{sec:appendix}

\subsection{Verification commands (F3 final wave)}\label{sec:appendix-f3}

Table~\ref{tab:appendix-f3-commands} reproduces the exact command sequence of
the F3 final verification wave
(\path{.omo/evidence/task-F3-tensixflow.log}, 2026-08-03, CPU-only box).
Every step exited 0. The wave additionally recorded: environment probes
confirming \texttt{/dev/tenstorrent} absent and no GPU present; emitted
project contents for both examples (host/reader/compute/writer.cpp +
CMakeLists.txt + README.md); regenerated report artifacts byte-identical to
the committed ones; captured PCC values for both targets on multiple seeds
(all $\ge 0.99998$, floor $0.9875$); and the offline HTML render check under
a network-blocked browser (zero requests, all six section anchors rendered).

\begin{longtable}{@{}p{0.7cm}p{10.6cm}p{3.4cm}@{}}
  \caption{Exact F3 final-wave command sequence with exit codes (all 0).}
  \label{tab:appendix-f3-commands}\\
  \toprule
  \textbf{\#} & \textbf{Exact command} & \textbf{Result (rc)} \\
  \midrule
  \endfirsthead
  \multicolumn{3}{@{}l}{\small\itshape Table \thetable{} (continued)}\\
  \toprule
  \textbf{\#} & \textbf{Exact command} & \textbf{Result (rc)} \\
  \midrule
  \endhead
  \bottomrule
  \endfoot
  1a &
  \begin{minipage}{\linewidth}\begin{lstlisting}
nix develop --command bash -c 'cmake -G Ninja -B build -DMLIR_DIR=$MLIR_DIR -DLLVM_DIR=$LLVM_DIR -DLLVM_EXTERNAL_LIT=$(dirname $(which lit))/.lit-wrapped'
\end{lstlisting}\end{minipage}
  & configure rc=0 (MLIR 20.1.8 pinned) \\
  1b &
  \begin{minipage}{\linewidth}\begin{lstlisting}
nix develop --command bash -c 'cmake --build build'
\end{lstlisting}\end{minipage}
  & build rc=0 \\
  2 &
  \begin{minipage}{\linewidth}\begin{lstlisting}
nix develop --command bash -c 'cmake --build build --target check-tensixflow'
\end{lstlisting}\end{minipage}
  & rc=0; 25 = 24 passed + 1 unsupported \\
  3 &
  \begin{minipage}{\linewidth}\begin{lstlisting}
nix develop --command bash -c 'cd build && ctest --output-on-failure'
\end{lstlisting}\end{minipage}
  & rc=0; 11/11 passed \\
  4 &
  \begin{minipage}{\linewidth}\begin{lstlisting}
uv run --with numpy --with pytest pytest python/tensixflow/tests/ -q
\end{lstlisting}\end{minipage}
  & rc=0; 89 passed \\
  5 &
  \begin{minipage}{\linewidth}\begin{lstlisting}
build/bin/tensixflow-opt --tensixflow-infer-tile-layout --tensixflow-fuse-rmsnorm-matmul --tensixflow-plan-noc-sharding --tensixflow-allocate-circular-buffers --tensixflow-lower-to-ttkernel --tensixflow-emit-metalium=output-dir=<dir> examples/rmsnorm_matmul.mlir
# same command for examples/mlp_block.mlir
\end{lstlisting}\end{minipage}
  & rc=0 both examples; 6 files emitted each \\
  6 &
  \begin{minipage}{\linewidth}\begin{lstlisting}
build/bin/tensixflow-opt --tensixflow-plan-noc-sharding --tensixflow-allocate-circular-buffers examples/rmsnorm_matmul.mlir | build/bin/tensixflow-translate --tensixflow-to-json > reports/rmsnorm_matmul_plan.json
build/bin/tensixflow-report --in reports/rmsnorm_matmul_plan.json --out reports/rmsnorm_matmul_report.html
\end{lstlisting}\end{minipage}
  & rc=0; HTML regenerated byte-identical \\
  7 &
  \begin{minipage}{\linewidth}\begin{lstlisting}
uv run --with numpy --with pytest pytest python/tensixflow/tests/test_correctness.py -v
\end{lstlisting}\end{minipage}
  & rc=0; 11/11; all PCC $\ge 0.9875$, trip-wires hold \\
  8 &
  \begin{minipage}{\linewidth}\begin{lstlisting}
nix develop --command bash -c 'bash scripts/hw_acceptance.sh'
\end{lstlisting}\end{minipage}
  & rc=0; no-device branch: dry-run + 4 generated kernels compile, legacy rejected \\
  9 &
  \begin{minipage}{\linewidth}\begin{lstlisting}
uv run --with playwright python f3_offline_html_check.py
\end{lstlisting}\end{minipage}
  & rc=0; zero network requests, 6/6 anchors rendered \\
\end{longtable}

Steps 1--4 of Table~\ref{tab:appendix-f3-commands} were re-executed for this
report with identical outcomes (lit 25 = 24 + 1 unsupported; ctest 11/11;
pytest 89 --- Section~\ref{sec:software-validation}).

\subsection{Glossary}\label{sec:appendix-glossary}

\begin{longtable}{@{}p{3.4cm}p{11.3cm}@{}}
  \caption{Glossary of TensixFlow terms.}
  \label{tab:appendix-glossary}\\
  \toprule
  \textbf{Term} & \textbf{Definition} \\
  \midrule
  \endfirsthead
  \multicolumn{2}{@{}l}{\small\itshape Table \thetable{} (continued)}\\
  \toprule
  \textbf{Term} & \textbf{Definition} \\
  \midrule
  \endhead
  \bottomrule
  \endfoot
  Circular buffer (CB) & A per-core L1 FIFO of tiles staging data between the reader, compute, and writer tile programs. Sized depth $\times$ \texttt{kTileBytes} (a depth-160 buffer occupies $160 \times 2048 = 327{,}680$ B); indexed per the tt-metal \texttt{tt::CBIndex} convention \texttt{c\_0}\ldots\texttt{c\_31} (default pipeline: in0=\texttt{c\_0}, in1=\texttt{c\_1}, out=\texttt{c\_16}); synchronized via \texttt{cb\_reserve\_back}/\texttt{cb\_push\_back} (producer) and \texttt{cb\_wait\_front}/\texttt{cb\_pop\_front} (consumer). \\
  NoC & Network-on-chip: the Wormhole grid's asynchronous interconnect. TensixFlow plans per-core shard placement and tile movement over it (\texttt{noc\_async\_read}/\texttt{noc\_async\_write}, plus multicast to alleviate DRAM bandwidth contention) and estimates the resulting DRAM traffic in the plan. \\
  PCC & Pearson correlation coefficient. The primary numerical metric comparing a result tensor against the fp32-exact NumPy golden; the acceptance floor is 0.9875 in both the simulator harness and the (gated) on-device branch. \\
  Schedule IR & The tile-level \texttt{tensixflow.sched.*} dialect: a \texttt{sched.kernel} holding reader/compute/writer tile programs built from \texttt{tile\_loop} nests plus CB, NoC, and compute ops. It is the shared contract the emitter prints and the functional simulator executes (Section~\ref{sec:software-architecture}). \\
  Shard & The slice of a tensor assigned to one core's L1 by \texttt{tensixflow-plan-noc-sharding}, placed under the modeled 960 KiB CB/shard budget and routed over the NoC. \\
  Tile & The $32\times32$-element unit of data: bf16 tiles are $32\times32\times2 = 2048$ B. Tiles are the unit of CB depth, NoC moves, and compute ops. \\
  Trip-wire & An empirically calibrated PCC threshold set \emph{above} the good case's observed PCC, so a regression trips the wire before it ever approaches the 0.9875 acceptance floor (\texttt{test\_correctness.py}). \\
  Stub headers & Signature-only mirrors of the tt-metal API surface (\texttt{kernels/handwritten\_baselines/tt\_metal\_stubs/}): correct names and signatures, every body empty, pinned to tt-metal \texttt{@df36cd6}. A compile contract only --- no behavior. \\
  Structural-only validation & Validating generated kernels without execution: golden FileCheck patterns against the current tt-metal API plus the \texttt{-fsyntax-only} compile check against the stubs. Never run on hardware; no hardware numerics claimed. \\
  Disjunctive acceptance & The hardware acceptance policy of Section~\ref{sec:hardware-acceptance}: (on-device run AND PCC $\ge 0.9875$) OR (validated no-device substitute). Wired as an explicit \texttt{if/else}, never \texttt{\&\&\ldots||}, so a present-but-failing device can never be masked by the passing no-device branch. \\
\end{longtable}

\subsection{Deviation register}

Documented deviations from the original plan --- notably the no-device
hardware posture (todo 55's on-device entry point \texttt{SKIPPED}: no
Tenstorrent device is owned) and the deliberately unsubmitted upstream
contribution (todo 60 \texttt{SKIPPED} per explicit user instruction) ---
are recorded with their rationale in the decisions section
(Section~\ref{sec:decisions}). Both are policy-sanctioned: each has a
validated substitute (the no-device branch, which passes; the local
implementation plus submission documentation) rather than a silent gap.
%% INTEGRATOR NOTE: sec:decisions is owned by lane WA; if the reference is
%% still undefined at intermediate compile time, it resolves once lane WA's
%% content lands.
